%---------------------------------------------------------%
%   CHAPTER 4: USING THE LATEX CLASS                       %
%---------------------------------------------------------%
\chapter{Using the \texttt{umpthesis.cls} Class}
\label{chap:latex}

This chapter shows how to write and compile a PSM report with the \texttt{umpthesis.cls} class. The examples follow the structure of \texttt{UMP Proposal 2018.tex} \citep{umpthesiscls}.

\section{Prerequisites}

You need a working TeX Live or MacTeX installation with \texttt{pdflatex} and \texttt{bibtex}. The class depends on standard packages (memoir, geometry, hyperref, setspace, graphicx, tabularx, caption, enumitem, tocloft, tikz, etc.), which are included in a normal TeX Live installation.

\section{Project Layout}

Keep the report in a single folder:

\begin{verbatim}
New Proposal 2018/
|-- umpthesis.cls            % the class file
|-- logo-ump-color.png       % UMP logo for declarations
|-- UMP Proposal 2018.tex    % main file
|-- bib/
|   `-- mybiblio.bib         % bibliography database
|-- figure/
|   `-- example.png          % figures
`-- tex/
    |-- abstract.tex         % English abstract
    |-- abstrak.tex          % Malay abstract
    |-- introduction.tex     % Chapter 1
    |-- literature.tex       % Chapter 2
    `-- methodologies.tex    % Chapter 3
\end{verbatim}

\section{Main File Skeleton}

Create the main file with the class, margins, packages, and document profile:

\begin{verbatim}
\documentclass[english]{umpthesis}

\geometry{top=25mm,right=25mm,left=35mm,bottom=25mm,includefoot}

\usepackage{natbib}     % author-year citations
\usepackage{booktabs}   % three-line tables
\usepackage{graphicx}   % figures (already loaded by the class)

\setcounter{secnumdepth}{3}
\setcounter{tocdepth}{3}
\pagenumbering{roman}
\end{verbatim}

\section{Document Profile Commands}

Set the author, school, and thesis information before \texttt{\string\begin\string{document\string}}:

\begin{verbatim}
\author{Your Full Name}
\authorDOB{01/01/1990}
\authorStudID{ABC12345}
\authorNRICPassport{000000-00-0000}

\umpFaculty{Faculty of Computer Systems \& Software Engineering}
\umpSession{2026/2027}
\umpDegree{Bachelor of Computer Science}
\umpField{Computer Science}

\title{Your Report Title}
\titleWrapOne{Your Report Title}
\titleWrapTwo{}
\titleWrapThree{}

\svisorMain{Supervisor Name}
\svisorMainPos{Senior Lecturer}
% \svisorSecond{...}        % optional co-supervisor
% \svisorSecondPos{...}

\submitYear{2026}
\submitMonth{November}
\end{verbatim}

\section{Front Matter Commands}

The preliminaries are assembled in this order:

\begin{verbatim}
\begin{document}

\umpCoverPage              % front cover

\frontmatter               % roman page numbers start here
\maketitle                 % title page

% \umpCopyright{0}{0}{1}   % Declaration of Thesis and Copyright
                            % {CONFIDENTIAL}{RESTRICTED}{OPEN ACCESS}
% \umpLibLetter{RESTRICTED}{2}{Reason 1}{Reason 2}{Reason 3}
% \umpSvDeclaration         % Supervisor's Declaration
% \umpStudDeclaration       % Student's Declaration

\input{tex/abstrak}        % Malay ABSTRAK first
\input{tex/abstract}       % then English ABSTRACT

\tableofcontents
\listoftables
\listoffigures

\mainmatter                % arabic page numbers start here
\end{verbatim}

For a real report, remove the comment marks from the declaration commands that apply. The copyright declaration takes three values, each \texttt{1} for selected and \texttt{0} for not selected.

\section{Main Matter Commands}

Write each chapter in its own file and \texttt{\string\input} it:

\begin{verbatim}
\mainmatter

\input{tex/introduction}
\input{tex/literature}
\input{tex/methodologies}
\end{verbatim}

Inside a chapter file, use the class heading commands:

\begin{verbatim}
\chapter{Introduction}
\section{Background}          % 1.1
\subsection{History}          % 1.1.1
\subsubsection{Detail}        % 1.1.1.1

\begin{enumerate}[label=(\roman*)]  % (i), (ii), (iii)
\item First objective
\item Second objective
\end{enumerate}
\end{verbatim}

Cross-reference with \texttt{\string\label} and \texttt{\string\ref}:

\begin{verbatim}
\section{Background}
\label{sec:background}
See Section \ref{sec:background}.
\end{verbatim}

\section{Tables, Figures, and Equations}

\subsection{Table}

Place the caption above the table and use \texttt{booktabs} for the required three-line style:

\begin{verbatim}
\begin{table}[htbp]
\caption{Comparison of existing systems}
\label{tab:comparison}
\begin{tabularx}{\textwidth}{@{} l l l @{}}
\toprule
Feature & System A & System B \\
\midrule
Accuracy & High & Medium \\
Speed & Fast & Slow \\
\bottomrule
\end{tabularx}
\end{table}
\end{verbatim}

\subsection{Figure}

Place the caption below the figure:

\begin{verbatim}
\begin{figure}[htbp]
\centering
\includegraphics[width=0.8\textwidth]{figure/method.png}
\caption{Research methodology}
\label{fig:method}
\end{figure}
\end{verbatim}

\subsection{Equation}

Use \texttt{equation} for a numbered, centred equation:

\begin{verbatim}
\begin{equation}
y = mx + c
\end{equation}
\end{verbatim}

\section{References and Bibliography}

Create \texttt{bib/mybiblio.bib} with BibTeX entries, then cite with \texttt{natbib}:

\begin{verbatim}
\citep{key}   % (Author, Year)
\citet{key}   % Author (Year)
\end{verbatim}

At the end of the main file:

\begin{verbatim}
\bibliographystyle{plainnat}
\bibliography{bib/mybiblio}
\end{verbatim}

\texttt{plainnat} is author--year and close to APA. For strict APA 6th, replace it with a dedicated APA style such as \texttt{biblatex-apa} (requires changing from BibTeX to Biber).

\section{Compiling the Report}

Run the following commands in the project folder:

\begin{verbatim}
pdflatex -interaction=nonstopmode -halt-on-error main.tex
bibtex main
pdflatex -interaction=nonstopmode -halt-on-error main.tex
pdflatex -interaction=nonstopmode -halt-on-error main.tex
\end{verbatim}

The extra \texttt{pdflatex} passes resolve citations, the table of contents, and cross-references. Check the log for undefined citations or references before submitting:

\begin{verbatim}
grep -i undefined main.log
\end{verbatim}

\section{Cleaning Generated Files}

After the final compile, remove the auxiliary files and keep the sources and PDF:

\begin{verbatim}
rm -f *.aux *.bbl *.blg *.log *.out *.toc *.lof *.lot
\end{verbatim}
